\chapter{Первый спринт в Logra}
\label{ch:logra}

В главе~\ref{ch:fpga} мы выяснили, что ПЛИС не выполняет программу, а
\emph{становится схемой}. Чтобы такую схему проектировать, нужен инструмент, и
учиться мы будем сразу на нём. Для ПЛИС путь от идеи до платы слишком длинный,
чтобы проходить его руками: нужно описать схему текстом, проверить её, увидеть,
во что она превратилась, и только потом загрузить в кристалл. Все эти шаги
выполняет одна программа --- \term{Logra}. Она открывается в браузере, ничего
не требует устанавливать и показывает каждый этап пути, который подробно
разобран в главе~\ref{ch:flow}. Эта глава~--- короткая экскурсия, чтобы вы
знали, где что лежит. Что именно мы строим, объяснят следующие главы, начиная с
главы~\ref{ch:signals}: там каждая тема сразу проверяется в Logra.

Внутри Logra работают открытые инструменты: Yosys превращает текст в схему,
nextpnr размещает её на кристалле, Project Apicula упаковывает битстрим, а
Icarus Verilog проверяет поведение. Среда лишь собирает их вместе и рисует то,
что они вычислили.

\begin{buildbox}[Соберите сначала]
Откройте Logra и нажмите «Создать проект». Выберите «На основе примера»,
найдите пример «Бегущий огонь» и откройте его. Ничего не меняя, нажмите
«Синтезировать» (Ctrl+Enter) в правом верхнем углу. Пока идёт сборка, посмотрите
на схему: в ней есть кольцевой регистр, шесть триггеров, единица в которых
передаётся по кругу (что это такое, объясняет глава~\ref{ch:regs}). Дальше в
главе мы пройдём по вкладкам и найдём каждую часть этой схемы.
\end{buildbox}

\section{Карта среды}

Слева в Logra находится колонка вкладок. Они сгруппированы в том порядке, в
каком мы проходим путь от идеи до платы: «Узнать», «Создать», «Проверить» и
свёрнутая по умолчанию группа «Настоящая плата». Вкладки с готовым результатом
отмечены цветной точкой, у вкладки «Выводы» справа число назначенных ножек.

\begin{table}[H]
\centering
\caption{Группы и вкладки Logra, главы книги, которые к ним относятся}
\small
\begin{tabularx}{\textwidth}{l l X l}
\toprule
\textbf{Группа} & \textbf{Вкладка} & \textbf{Что в ней делают} & \textbf{Глава} \\
\midrule
Узнать & Обучение & Пошаговые уроки с вопросами и проверкой решения (курс «Твоя мастерская», 31 урок) & --- \\
\midrule
Создать & Код & Пишут описание схемы на Verilog & \ref{ch:verilog} \\
 & Конструктор & Собирают схему мышью из логических элементов, Verilog пишется сам & \ref{ch:logra} \\
 & Выводы & Назначают порты схемы на ножки корпуса & \ref{ch:board} \\
\midrule
Проверить & RTL & Смотрят схему из логических элементов, которую построил синтезатор & \ref{ch:bridge} \\
 & Диаграммы & Подают сигналы и смотрят временные диаграммы & \ref{ch:sim} \\
 & Лаборатория & Подключают к плате детали и запускают схему «как на столе» & \ref{ch:examples} \\
\midrule
Настоящая плата & Gowin & Видят LUT, триггеры и блоки сложения после синтеза & \ref{ch:inside} \\
 & Кристалл & Размещение и трассировка, отчёт о частоте, битстрим & \ref{ch:flow} \\
 & Подключение & Загрузка в плату по USB, UART-терминал & \ref{ch:flow} \\
\bottomrule
\end{tabularx}
\end{table}

В верхней строке справа находятся кнопка «Создать проект», переключатель
языка (RU/EN), переключатель темы, справка и главная кнопка
\textbf{«Синтезировать»} (Ctrl+Enter), которая запускает весь маршрут. В
нижней строке состояния написано «Yosys готов» (это название программы
Yosys, а не среды) и «Icarus Verilog готов».

\begin{diybox}[Карта среды]
\textbf{Задание.} Пройдите по всем вкладкам и убедитесь, что вы понимаете
группы.\\
\textbf{В Logra.} Откройте проект «Бегущий огонь», нажмите «Синтезировать».\\
\textbf{Ожидаемый результат.} У вкладок RTL, Gowin и Кристалл появились
цветные точки: результат готов; у «Выводов» стоит число назначенных ножек.\\
\textbf{Проверка.} Разверните группу «Настоящая плата» (стрелка слева от
названия) и найдите в ней три вкладки. Запишите, в какой группе лежит каждая из
вкладок Код, Диаграммы, Кристалл.\\
\stars{1} Повторите: найдите кнопку смены языка и верните интерфейс на русский.
\end{diybox}

\section{Новый проект}

Работа начинается с кнопки «Создать проект» в верхней строке. Окно
предлагает пять способов начать: «С нуля» (написать Verilog самому),
«Из логических элементов» (собрать схему мышью), «На основе примера» (в
каталоге 275 готовых проектов), «Импорт проекта» (из файла) и «Папка проекта
Gowin». Верхние кнопки окна: «Создание», «Мои проекты», «Текущий проект»
и «Синхронизация». Мы выберем «С нуля»: нужно только дать проекту имя и
указать имя верхнего модуля, обычно его называют \code{top}.

\shot{app-wizard}{Окно создания проекта: пять способов начать}{fig:y-wizard}

Вариант «Из логических элементов» открывает \term{конструктор}: схему собирают
мышью из библиотеки, а Logra сам пишет Verilog. Это тот же путь, что на
бумаге (так мы будем делать в главах~\ref{ch:gates}--\ref{ch:arith}), только с
мгновенной проверкой. Библиотека устроена по
образцу графического редактора Quartus: вентили с тремя, четырьмя и восемью
входами, триггеры (D, T, JK, RS), защёлка, регистры и счётчики, ОЗУ и ПЗУ,
мультиплексоры, дешифраторы и шифраторы, сумматоры, компаратор, умножитель и
делитель, модуль числа, код Грея и дешифратор для семисегментного индикатора.
Готовые наборы названы как привычные микросхемы серии 74: 74138, 74147,
74148, 7442, 7447, 74184, 7483, 7485. Кнопка «Предпросмотр» позволяет сразу
проверить схему: щёлкать по входам и подавать такт.

\begin{diybox}[Первый проект]
\textbf{Задание.} Создайте проект и соберите в нём вентиль «исключающее ИЛИ» из
библиотеки.\\
\textbf{В Logra.} «Создать проект» $\to$ «Из логических элементов»; поставьте
два входа, элемент XOR и выход.\\
\textbf{Ожидаемый результат.} В «Предпросмотре» выход равен 1 ровно тогда, когда
входы различны (таблица истинности из главы о вентилях).\\
\textbf{Проверка.} Откройте вкладку «Код» и найдите строку \code{assign} с
оператором \code{\^{}}; нажмите «Синтезировать» и убедитесь, что в RTL остался
один элемент XOR.\\
\stars{2} Измените схему так, чтобы получился XNOR, и сравните с таблицей.
\end{diybox}

\section{Код}

Вкладка «Код» --- это редактор с подсветкой Verilog. Ошибки подчёркиваются
по мере набора, а сохранение происходит само (внизу написано «Автосохранено»).
Под редактором есть «Журнал сборки» и «Диагностика». Для разбора возьмём проект
«Бегущий огонь»: единица в шестиразрядном кольцевом регистре перемещается на
следующий светодиод восемь раз в секунду. Это схема из главы~\ref{ch:examples},
но порты названы так, как их называет плата: \code{clk}, \code{reset} и
\code{led\_n}.

\shot{app-code}{Редактор кода. Кнопка «Синтезировать (Ctrl+Enter)» в углу запускает весь маршрут}{fig:y-code}

\begin{infobox}[Почему \code{led\_n}]
Суффикс \code{\_n} --- обычное соглашение: сигнал с активным \emph{низким}
уровнем. Светодиоды платы горят, когда на выводе 0, поэтому в схеме мы
подаём на них инвертированное значение кольца.
\end{infobox}

\begin{diybox}[Свой параметр]
\textbf{Задание.} Замедлите бегущий огонь вдвое.\\
\textbf{В Logra.} Вкладка «Код», строка \code{localparam DIV = 27\_000\_000 / 8;}.\\
\textbf{Ожидаемый результат.} После замены \code{8} на \code{4} огонь переходит
четыре раза в секунду.\\
\textbf{Проверка.} Нажмите «Синтезировать» и откройте «Лабораторию»: полоса из
шести светодиодов должна переключаться вдвое реже. Если «Диагностика» показывает
ненулевое число, прочитайте сообщение и найдите опечатку.\\
\stars{2} Сделайте так, чтобы огонь бежал в обратную сторону (поменяйте
порядок битов в склейке).
\end{diybox}

\section{RTL: схема, которую построил синтезатор}

Нажмём «Синтезировать». Меньше чем за секунду Yosys читает текст и строит
схему. Вкладка RTL (register-transfer level, «уровень регистровых передач»:
описание схемы как набора регистров и логики между ними) показывает её так,
как рисуют на бумаге: прямоугольники для регистров, условные знаки для
вентилей, кружки с плюсом для сумматоров. Это ровно то, что мы будем
изучать в главах про логические элементы, сумматоры и триггеры, только
нарисованное автоматически.

\shot{br-adder-rtl}{RTL-схема на примере полного сумматора (глава~\ref{ch:bridge}): вентили XOR, AND и OR, подписи с номерами строк кода и справа инспектор}{fig:y-rtl}

Щелчок по элементу открывает справа «Инспектор»: входы, выходы, параметры и
строку исходного текста, из которой элемент получился. Двойной щелчок
переносит в редактор. Колесо мыши меняет масштаб, перетаскивание двигает схему.

\begin{ideabox}
Среда показывает схему \emph{до} привязки к конкретному кристаллу. Это
«идеальная» схема из вентилей и триггеров. Внизу панели написано: «До
отображения на FPGA». Что из неё получится на настоящем кристалле, видно на
вкладке Gowin.
\end{ideabox}

\begin{diybox}[Прочитайте RTL]
\textbf{Задание.} Найдите на схеме «Бегущего огня» кольцо из шести триггеров и
сумматор счётчика.\\
\textbf{В Logra.} Вкладка RTL после «Синтезировать».\\
\textbf{Ожидаемый результат.} Видны регистры, элемент сложения (кружок с
плюсом), сравнение \code{cnt == DIV - 1} и инверторы перед выходом \code{led\_n}.\\
\textbf{Проверка.} Щёлкните по сумматору: в «Инспекторе» должна быть указана
строка кода \code{cnt + 1}; двойной щелчок переносит в редактор на эту строку.
\end{diybox}

\section{Выводы: от порта к ножке}

Порт \code{clk} в Verilog --- просто имя. Чтобы он заработал, нужно сказать,
к какой ножке корпуса подключён генератор 27~МГц. Для этого служит вкладка
«Выводы» с подвкладками «Планировщик выводов» и «Текст .cst». Слева таблица
портов вашей схемы, справа корпус GW2AR-18 QN88 с подписями всех ножек. Подсказка
из трёх шагов: выберите сигнал, выберите контакт, нажмите «Назначить».
Кнопка «Назначить LED по порядку» расставляет шесть светодиодов сама. Среда
запишет строки \code{IO\_LOC} в файл \code{board.cst} --- тот самый файл
ограничений, который мы разбирали в разделе~\ref{sec:cst}.

\shot{app-pins}{Планировщик выводов. Для портов с подходящими именами ножки подставлены сами}{fig:y-pins}

\begin{diybox}[Найдите ножку]
\textbf{Задание.} Определите, к какой ножке ПЛИС подключён вывод \code{clk} и
какую строку Logra записал в файл ограничений.\\
\textbf{В Logra.} Вкладка «Выводы»: подвкладки «Планировщик выводов» и «Текст .cst».\\
\textbf{Ожидаемый результат.} \code{clk} на ножке 4; светодиоды \code{led\_n[0..5]}
на ножках 15--20 (табл.~\ref{tab:pins}).\\
\textbf{Проверка.} Откройте «Текст .cst»: строки \code{IO\_LOC "clk" 4;} и
\code{IO\_LOC "led\_n[0]" 15;} должны совпасть с планировщиком.\\
\stars{2} Намеренно испортите имя порта в \code{.cst} и посмотрите, какое
сообщение даст сборка (см. предупреждение в разделе~\ref{sec:cst}).
\end{diybox}

\section{Gowin: во что превратилась схема}

Вкладка Gowin показывает результат следующего шага --- отображения на
примитивы кристалла. Вместо абстрактных вентилей здесь LUT, триггеры DFF,
блоки арифметики ALU и буферы ввода-вывода. Если щёлкнуть по LUT, «Инспектор»
выведет блок «Таблица LUT на 3 входа», таблицу конфигурации INIT и сворачиваемую
«Таблицу истинности». Именно так наглядно видно утверждение из
главы~\ref{ch:inside}: логический элемент в ПЛИС --- это таблица в памяти.

\shot{br-adder-lut}{Тот же сумматор после отображения на примитивы Gowin: вместо вентилей остались две LUT, справа в инспекторе видна таблица LUT и значение INIT}{fig:y-gowin}

\begin{diybox}[Найдите LUT]
\textbf{Задание.} Найдите в схеме одну LUT и сравните её таблицу с функцией из
кода.\\
\textbf{В Logra.} Откройте пример «Полный сумматор», нажмите «Синтезировать»,
вкладка Gowin; щёлкните по LUT.\\
\textbf{Ожидаемый результат.} В «Инспекторе» таблица на 3 входа с
конфигурацией INIT: у суммы \code{s = a \^{} b \^{} cin} единицы стоят в строках
с нечётным числом единиц на входах.\\
\textbf{Проверка.} Сверьте таблицу с таблицей истинности сумматора (глава~\ref{ch:arith}) и
перечислите строки, где выход равен 1.\\
\stars{2} Определите по значению INIT, какая функция записана в LUT переноса.
\end{diybox}

\section{Диаграммы: поведение во времени}

Схема нарисована, но работает ли она? Вкладка «Диаграммы» запускает
симулятор Icarus Verilog и рисует временные диаграммы сигналов
(рис.~\ref{fig:y-wave}). Для комбинационных схем есть сценарий «Перебрать
входы»: Logra сам подаст все комбинации входных сигналов по очереди. Для
последовательных схем среда сама вырабатывает такт и короткий импульс сброса.
Подробнее о симуляции --- в главе~\ref{ch:sim}.

\shot{br-adder-wave}{Все восемь комбинаций входов полного сумматора: так проверяют таблицу истинности}{fig:y-wave}

\subsection{Сценарий: события, циклы и выбор сигналов}

Для схем с памятью одного перебора мало: нужно подать входы в нужные
моменты. Кнопка «Настроить сценарий» открывает панель с четырьмя вкладками.
\emph{Запуск} задаёт тактовый вход, сигнал сброса, период и число циклов.
\emph{Входы} хранит постоянные значения. На вкладке \emph{События} можно
добавлять шаги: «в периоде 5 вход \code{c5} равен 1, в периоде 6 --- снова 0».
Вкладка \emph{Проверки} позволяет записать ожидаемые значения, и Logra
сравнит их с результатом симуляции.

Если сигналов много, меню «Сигналы» над диаграммой позволяет скрыть
служебные цепи и оставить только нужные, а кнопка «Показать весь интервал»
вписывает всю симуляцию в окно. Так построены диаграммы автомата и очереди
в главе~\ref{ch:fsmdesign}: все входные импульсы заданы на вкладке «События».

\begin{diybox}[Проверка таблицы истинности]
\textbf{Задание.} Проверьте полный сумматор перебором всех входов.\\
\textbf{В Logra.} Пример «Полный сумматор», вкладка «Диаграммы», «Перебрать
входы».\\
\textbf{Ожидаемый результат.} Восемь интервалов: пара \code{cout, s} равна
двоичной сумме \code{a + b + cin}.\\
\textbf{Проверка.} Выберите любые три интервала и вычислите сумму в уме;
сравните с диаграммой.\\
\stars{2} Для «Бегущего огня» откройте «Настроить сценарий» и увеличьте
число циклов; найдите момент, когда единица в кольце вернулась на первый
светодиод.
\end{diybox}

\section{Кристалл: размещение и битстрим}

Кнопка «Разместить и собрать» на панели «Размещение и трассировка» запускает
nextpnr. Он решает, на какую именно ячейку кристалла посадить каждую LUT и
каждый триггер, и проводит между ними соединения. На «Карте размещения на
кристалле» цветными квадратами видно, какая часть кристалла занята. Над картой
счётчики ресурсов: для «Бегущего огня» --- 15 LUT, 29 триггеров, 22 ALU и 8 ножек
ввода-вывода. Рядом отчёт о частоте: nextpnr показывает 338,87~МГц
при требуемых 27~МГц, то есть запас примерно в двенадцать раз.

\shot[0.82\textwidth]{app-chip}{Карта размещения на кристалле и результат nextpnr: 338,87 МГц при 27 МГц}{fig:y-chip}

После успешной сборки доступны файлы \code{.fs} (битстрим), \code{JSON} и журнал.

\begin{diybox}[Запас по частоте]
\textbf{Задание.} Найдите, сколько ресурсов занимает «Бегущий огонь» и какая у
него максимальная частота.\\
\textbf{В Logra.} Группа «Настоящая плата», вкладка «Кристалл», «Разместить и
собрать».\\
\textbf{Ожидаемый результат.} Счётчики LUT, FF, ALU, IO и строка вида
«338.87 / 27.00 МГц. Частота достигнута».\\
\textbf{Проверка.} Разделите максимальную частоту на 27 и запишите запас;
затем увеличьте в поле «Частота, МГц» требование до 400 и пересоберите: что
сообщает отчёт?\\
\stars{2} Сравните число FF с числом разрядов в коде (22 в счётчике, 1 для
\code{tick}, 6 в кольце).
\end{diybox}

\section{Лаборатория}

Прежде чем грузить схему в настоящую плату, её можно запустить на
виртуальном столе. Лаборатория показывает плату Tang Nano 20K на макетной
плате. Сверху кнопки «Пауза», «Такт», «Сброс», выбор скорости, кнопки платы
S1 и S2 и счётчик «Время на плате». К плате можно подключить детали:
светодиоды, индикаторы, датчики, экраны, память. Схема исполняется в браузере
непрерывно, а детали ведут себя как настоящие, включая последовательные
протоколы I$^2$C, SPI и UART (глава~\ref{ch:uart}).

\shot[0.82\textwidth]{app-lab}{Бегущий огонь на виртуальной плате. Слева список из 85 деталей}{fig:y-lab}

\begin{diybox}[Виртуальный стол]
\textbf{Задание.} Поставьте на стол светодиод и подключите его к выходу схемы.\\
\textbf{В Logra.} Вкладка «Лаборатория»: деталь «Светодиод» из группы
«Светодиоды», провод от контакта платы.\\
\textbf{Ожидаемый результат.} Светодиод мигает в такт схеме; кнопки S1 и S2
меняют состояние входов.\\
\textbf{Проверка.} Нажмите «Пауза» и «Такт»: схема делает шаг за шагом, и
вы видите, когда меняется состояние. Найдите в поиске деталей слово
«светофор».\\
\stars{2} Найдите деталь «Светофор» и подключите её к схеме светофора из примеров.
\end{diybox}

\section{Подключение и обучение}

Вкладка «Подключение» загружает битстрим в плату прямо из браузера по
WebUSB. Переключатели вверху: «Загрузить схему», «UART-терминал» и «Отладчик».
Три шага: подключить плату, выбрать файл, нажать «Загрузить в SRAM». В
режиме «UART-терминал» в окно приходят символы, отправленные схемой по
последовательной линии, и из него можно отправить символы в схему.

\shot[0.82\textwidth]{app-conn}{Загрузка битстрима в плату}{fig:y-conn}

А вкладка «Обучение» --- самостоятельный курс «Твоя мастерская» из тридцати
одного урока для тех, кто видит ПЛИС впервые: от «что такое бит» до
собственных схем с проверкой решений симулятором. Есть режимы «С наставником»
и «Свободные задания».

\shot[0.82\textwidth]{app-learn}{Курс «Твоя мастерская»: уроки с проверкой}{fig:y-learn}

\begin{diybox}[Первый урок и загрузка]
\textbf{Задание.} Пройдите первый урок курса и найдите шаги загрузки в плату.\\
\textbf{В Logra.} Вкладка «Обучение»: урок «Что у нас в руках»; затем вкладка
«Подключение».\\
\textbf{Ожидаемый результат.} Урок отмечен выполненным; на «Подключении»
видны три шага и кнопка «Загрузить в SRAM» (без платы она остаётся
недоступной).\\
\textbf{Проверка.} Если плата есть, загрузите схему и убедитесь, что огонь
бежит по настоящим светодиодам; без платы сверьтесь с «Лабораторией».
\end{diybox}

\section{Итог спринта}

\begin{taskbox}[Задание]
Создайте в Logra проект с нуля, назовите модуль \code{top} и введите
описание с рис.~\ref{fig:y-code}. Нажмите «Синтезировать» и последовательно
пройдите вкладки RTL, Gowin, Кристалл, Лаборатория. На каждой найдите:
сколько триггеров в схеме (29: шесть в кольце, двадцать два в делителе и один для сигнала tick),
какую максимальную частоту обещает nextpnr и в каком порядке загораются
светодиоды.
\end{taskbox}

\begin{diybox}[Спринт «Первый путь»]
\textbf{Результат спринта.} Проект с модулем \code{top}, прошедший синтез,
размещение и запуск в Лаборатории, и ваша таблица из четырёх чисел
(триггеры, LUT, частота, номер ножки \code{clk}).\\
\stars{3} \textbf{Этап проекта «Охранная панель»:} создайте проект с этим
именем и сохраните в нём первый вариант (хотя бы один вход-датчик и один
светодиод-индикатор). Критерий приёмки: синтез проходит без ошибок, в
«Выводах» назначены все порты, в «Лаборатории» светодиод реагирует на кнопку S1.
\end{diybox}

\section{Что дальше}

Теперь у вас есть рабочая среда и первый проект, но вы пока только нажимали
кнопки: почему схема из «Бегущего огня» устроена именно так, какие сигналы в ней
бегут и как превратить событие в мире (дверь открылась) в ноль или единицу? Это
вопрос следующей главы~\ref{ch:signals}: мы начнём с сигнала и временной
диаграммы, а проект «Охранная панель», созданный здесь, будет расти вместе с
книгой.
